\section{Evaluation Protocol and Results}
\label{sec:eval}

The evaluation protocol is designed to detect three failure modes simultaneously: (a) raw factual ignorance, (b) self-preference bias from generator overlap, and (c) inflation due to questions solvable from world knowledge alone.

\subsection{Splits}
We define four evaluation slices over the OenoBench corpus:
\begin{itemize}\itemsep -1pt
  \item \textbf{Full corpus} (\PLACEHOLDER{$N=5{,}000$}) --- the headline benchmark.
  \item \textbf{Per-generator held-out splits} --- for each generator model $g$, the slice of questions \emph{not} produced by $g$. Used to compute self-vs-other gaps.
  \item \textbf{Closed-book reserve} (\texttt{cb\_reserve}) --- questions flagged by agent B2 as world-knowledge solvable. Reported separately as a pure-recall slice.
  \item \textbf{Expert control set} ($N=300$) --- the WSET-Diploma--authored, zero-LLM control. Functions as a ground-truth comparator against which we calibrate the rest of the corpus.
\end{itemize}

\subsection{Self-Preference Score}
For each model $m$, we define
\[
  \mathrm{SPS}(m) \;=\; \mathrm{Acc}(m \mid \text{questions generated by } m) \;-\; \mathrm{Acc}(m \mid \text{questions generated by other models}).
\]
A pipeline that successfully neutralises generator bias should produce $|\mathrm{SPS}(m)| < 3\,\mathrm{pp}$ for every $m$ in our model set. Larger gaps imply either (i) the generator's own questions are systematically easier for it than for peers, or (ii) the audit pipeline is failing to remove stylistic fingerprints. We treat $|\mathrm{SPS}|>3\,\mathrm{pp}$ as a corpus-level fail and trigger a re-paraphrase pass on the affected generator's slice.

\subsection{Models evaluated}
We evaluate \PLACEHOLDER{final list of frontier LLMs}: Claude (Opus 4.6, Sonnet 4.6), GPT (GPT-4o, GPT-5), Gemini (1.5 Pro, 3.1 Pro), Llama (3.1 405B), and at least one open small model as a baseline. Each model is run zero-shot, then with chain-of-thought prompting; we report the better of the two as the headline number, and both separately in the appendix.

\subsection{Headline results}
\PLACEHOLDER{Table 3: full-corpus accuracy per model, broken down by domain pillar (regions / varieties / viticulture / winemaking / producers / business). 8 rows, 7 columns. Include 95\% bootstrap CI.}

\PLACEHOLDER{Table 4: Self-Preference Score per model on the per-generator held-out slice. Pass threshold $|\mathrm{SPS}|<3\,\mathrm{pp}$.}

\PLACEHOLDER{Table 5: closed-book reserve vs.\ active-corpus accuracy gap per model. The closed-book reserve is, by construction, easier; we expect a +X\,pp gap, and use deviations from this expectation to flag overfitting to the rest of the corpus.}

\PLACEHOLDER{Figure 1: per-domain accuracy radar chart, all frontier models overlaid.}

\subsection{Comparison to expert control set}
\PLACEHOLDER{Table 6: model accuracy on the 300-question expert control set vs.\ the LLM-generated bulk of the corpus. A model that scores significantly differently on the two slices reveals that LLM-generated questions are either easier (generation-side blind-spots) or harder (artefacts) than expert-written ones; the gap is itself a benchmark-validity diagnostic.}

\subsection{Error analysis}
\PLACEHOLDER{Qualitative error analysis: 200 randomly-sampled wrong answers, hand-categorised by error type (factual confusion, regional misattribution, vintage error, regulatory misreading, sensory descriptor mismatch). Discussion of which domains and which generation strategies surface which error classes most often.}
